\chapter{Related Work}

\section{Text-Based Legal Judgment Prediction}

Early LJP relied on bag-of-words and TF--IDF features fed into linear or tree-based classifiers. Aletras et al.\ predicted ECHR violation judgments with around 79\% accuracy from n-gram and topic features over case text, with the facts section emerging as the strongest signal~\cite{aletras2016echr}, and Katz et al.\ predicted US Supreme Court outcomes across nearly two centuries with a time-evolving random forest over institutional metadata and voting patterns~\cite{katz2017supreme}. Medvedeva et al.\ then showed that predicting genuinely future cases from past ones substantially lowers performance, exposing temporal leakage as a structural problem in the literature~\cite{medvedeva2020echr}, and a subsequent methodological audit found that most LJP publications do not perform genuine prospective prediction at all~\cite{medvedeva2023doitright}. The present thesis treats leakage control as a first-class concern for that reason.

The neural wave replaced manual features with learned representations. LEGAL-BERT and the LexGLUE benchmark established that domain-adapted transformers consistently outperform generic ones across legal classification tasks~\cite{chalkidis2020legalbert, chalkidis2022lexglue}. The persistent weakness of standard transformers on legal text is document length: most judgments exceed 512 tokens and even extended-context architectures incur substantial cost, with hierarchical TF--IDF variants remaining competitive on long inputs~\cite{mamakas2022longlegal}. Both feature-engineered and transformer pipelines also share a more fundamental limit: a flat or hierarchical document representation cannot exchange information across cases that share legally meaningful objects. The present thesis treats text embeddings as node features within a relational representation rather than as the complete model.

\section{Graph-Based Approaches}

Graph-based methods were developed precisely to model the relational structure that flat text pipelines miss. TopJudge formalised dependencies among legal subtasks, law articles, charges, and penalties, as a directed acyclic graph, showing that joint prediction over structured outputs improves performance compared to treating each subtask independently~\cite{zhong2018topjudge}. Subsequent work expanded this relational view by incorporating retrieved precedents and law-article graphs, with more recent systems using LLM-and-retrieval pipelines to bring external legal scaffolding into prediction~\cite{wu2023pljp, santosh2024precedents}.

Within Indian LJP, graph-based work is growing. Khatri et al.\ explored GNNs for Indian legal judgment prediction and raised important questions about the role of case-graph design and demographic shortcuts~\cite{khatri2023indianGNN}. More recent large-scale Indian work has extended coverage and explored RAG-style pipelines~\cite{nigam2025nyayaanumana}. Existing graph approaches nonetheless leave gaps that this thesis addresses: many are built for civil-law charge-prediction datasets, rely on a single court, or use graph structure to improve prediction without comparable attention to decision-text or identity leakage. The present work is closest in spirit to this line but narrows its claim: the graph should improve prediction by exposing legal structure across cases while explicitly blocking shortcuts that rely on post-decision information.

\section{Heterogeneous GNNs and the Motivation for HGT}

A homogeneous GNN is too blunt for legal graphs: a case node, a statute node, and a precedent node play different epistemic roles, and edges such as \emph{cites\_precedent} and \emph{heard\_in} should not be collapsed into a single message-passing rule. Among heterogeneous alternatives, HGT~\cite{hu2020hgt} fits the design requirements of this thesis best: its node-type-dependent projections and edge-type-dependent attention let the model weight heterogeneous neighbours without the relation-specific weight matrices of R-GCN~\cite{schlichtkrull2018rgcn} growing coarse as the type count rises, and without the handcrafted meta-path templates required by HAN~\cite{wang2019han}. The graph in this thesis has many node and edge types and a deliberately cross-domain structure, which is exactly the regime HGT was designed for.

\section{Indian Legal NLP and OpenNyAI}

Indian legal NLP developed later than English-language benchmarks for US or European law, in part due to data and tooling scarcity. Early work on rhetorical role identification showed that Indian judgments have a recoverable discourse structure, preamble, facts, arguments, ratio, and ruling, that is useful for downstream tasks but not explicitly marked in raw text~\cite{bhattacharya2019rhetoricalroles}. This made legal document structuring a first-class problem rather than a preprocessing convenience.

OpenNyAI is the most practically important development for this thesis, bringing together legal named entity recognition, rhetorical-role prediction, and summarisation for Indian court judgments into an open, modular pipeline~\cite{kalamkar2022ner, opennyai2024docs, opennyai2024about}. On the benchmarking side, ILDC introduced a large corpus of Indian Supreme Court cases for judgment prediction and explanation~\cite{malik2021ildc}, while IL-TUR broadened the scope to a multi-task legal understanding and reasoning benchmark~\cite{joshi2024iltur}. Newer rhetorical-role resources such as LegalSeg confirm that structural modelling continues to matter for Indian judgments~\cite{nigam2025legalseg}.

These contributions are substantial, but they do not resolve the core challenge addressed here: how to represent Indian cases so that a model can exploit legal structure across cases without quietly learning identity or post-decision shortcuts. This thesis builds on the Indian legal NLP ecosystem rather than apart from it, combining OpenNyAI's pipeline components with a leakage-audited heterogeneous graph.

\section{Gaps This Work Addresses}

The literature reviewed above shows that legal judgment prediction has advanced substantially, but along partly disconnected lines. Classical models demonstrated that outcomes are statistically predictable; transformer models improved semantic representation; graph methods showed the value of relational structure; Indian NLP resources made domain-specific preprocessing feasible. Yet most prior work sacrifices at least one of the following: \emph{scale}, \emph{cross-domain coverage}, \emph{cross-case legal structure}, or \emph{leakage-aware experimental design}.

What remains comparatively rare is an Indian LJP system that operates over a large multi-domain corpus, consolidates multiple hearings per case, extracts legal entities and rhetorical structure, builds an explicitly heterogeneous cross-case graph, and evaluates under a design intended to resist both decision-text leakage and identity-based shortcuts. That is the niche this thesis occupies. Its claim is narrow and, for that reason, defensible: when Indian court judgments are converted into a leakage-audited heterogeneous graph, cross-case legal structure becomes usable in ways that flat text pipelines cannot easily match.